% This is a sample LaTeX input file.  (Version of 12 August 2004.)
%
% A '%' character causes TeX to ignore all remaining text on the line,
% and is used for comments like this one.


\documentclass[twocolumn]{article}      % Specifies the document class
\usepackage{authblk}
\usepackage{amsmath}
\usepackage{algorithm}
\usepackage{algpseudocode}
\usepackage{graphicx}
\usepackage{ragged2e}
\makeatletter
\def\BState{\State\hskip-\ALG@thistlm}
\makeatother
                             % The preamble begins here.
\title{The Largest and the Smallest Pythagorean Number from randomly generated 1000 Natural Numbers}  % Declares the document's title.
 \author{Akshay Gupta~{IIT2017505}, \hspace{6pt}Naman Deept~{IIT2017507}  \hspace{6pt}Snigdha Dobhal~{IIT2017506}
 }
\date{\today}      % Deleting this command produces today's date.
\renewcommand\Authands{ and }
\newcommand{\ip}[2]{(#1, #2)}
                             % Defines \ip{arg1}{arg2} to mean
                             % (arg1, arg2).

%\newcommand{\ip}[2]{\langle #1 | #2\rangle}
                             % This is an alternative definition of
                             % \ip that is commented out.

\begin{document}             % End of preamble and beginning of text.
\begin{titlepage}

\newcommand{\HRule}{\rule{\linewidth}{0.5mm}} % Defines a new command for the horizontal lines, change thickness here

\center % Center everything on the page
 
%----------------------------------------------------------------------------------------
%   HEADING SECTIONS
%----------------------------------------------------------------------------------------

\textsc{\LARGE Assignment 1}\\[1.5cm] % Name of your university/college
\includegraphics[scale=0.7]{iiitlogo.jpg}\\[1cm] % Include a department/university logo - this will require the graphicx package
\textsc{\Large Design and Analysis of Algorithm}\\[0.5cm] % Major heading such as course name
\textsc{\large IDAA432C}\\[0.5cm] % Minor heading such as course title

%----------------------------------------------------------------------------------------
%   TITLE SECTION
%----------------------------------------------------------------------------------------

\HRule \\[0.4cm]
{ \large \bfseries Largest and Smallest Pythagorean Numbers from a Randomly Generated Array of Natural Numbers. }\\[0.4cm] % Title of your document
\HRule \\[1.5cm]
 
%----------------------------------------------------------------------------------------
%   AUTHOR SECTION
%----------------------------------------------------------------------------------------

\begin{minipage}{0.4\textwidth}
\begin{flushleft} \large
\textbf{\emph{Submitted By:}}\\
Akshay Gupta \textbf{(IIT2017505)}\\
Naman Deept \textbf{(IIT2017507)}\\
Snigdha Dobhal \textbf{(IIT2017506)}
\end{flushleft}

\end{minipage}\\[2cm]

% If you don't want a supervisor, uncomment the two lines below and remove the section above
%\Large \emph{Author:}\\
%John \textsc{Smith}\\[3cm] % Your name

%----------------------------------------------------------------------------------------
%   DATE SECTION
%----------------------------------------------------------------------------------------

{\large \today}\\[2cm] % Date, change the \today to a set date if you want to be precise

\vfill % Fill the rest of the page with whitespace

\end{titlepage}
\maketitle                   % Produces the title.


    % Produces section heading.  Lower-level
                             % sections are begun with similar 
                             % \subsection and \subsubsection commands.

\section{Abstract} 
\textbf{T}his paper introduces several algorithms to find the largest and the smallest \textbf{Pythagorean} Numbers from the randomly generated array of \textbf{N}atural Numbers and then the most optimal approach to the problem is considered to efficiently solve the problem.\\
\textbf{Keywords}: \textit{Pythagorean}, \textit{Pythagorean Triples}, \textit{Primes}, \textit{Euclid's Formula}, \textit{Pythagorean prime}
\section{Introduction}
{ \textbf{Pythagorean Number:}} A Natural number \textbf{a} is said to be a Pythagorean number if it can be written as sum of the square of two smaller Natural numbers \textbf{b} and \textbf{c} , such that ,
\begin{center}\textbf{$a^2 = b^2 + c^2$}\end{center}
Such a triple is commonly written \textbf{(a, b, c)}, and a well-known example is (3, 4, 5). If (a, b, c) is a Pythagorean triple, then so is \textbf{(ka, kb, kc)} for any positive integer k. A primitive Pythagorean triple is one in which a, b and c are co-prime (that is, they have no common divisor larger than 1). A triangle whose sides form a Pythagorean triple is called a Pythagorean triangle, and is necessarily a right triangle. Any Pythagorean Number is a multiple of a \textbf{Pythagorean} prime \textit{(n)} which is prime as well as of the form of \textit{4n+1}.\\\ When searching for integer solutions, the equation \textbf{$a^2 + b^2 = c^2$} is a Diophantine equation. Thus Pythagorean triples are among the oldest known solutions of a nonlinear Diophantine equation.
\section{Proposed Method}
\textbf{Input}: Given 1000 randomly generated Natural numbers in an array , the aim is to find the largest and the smallest Pythagorean Number from this array of natural numbers ranging in between $10^4$ to $10^6$.
\subsection {Naive approach}
First method could be to check whether there exist any number \textbf{i} less than  the given number \textbf{n} such that $n^2 - i^2$ is a perfect square . If such a number exist then the given number is a Pythagorean Number . 
\subsubsection{Algorithm for the Naive Approach}
\begin{algorithm}
\begin{algorithmic}[1]
\Procedure{checkPythagorean}{$n$}
\For {$i \gets 3 \hspace{4pt} to \hspace{4pt} {n} $}
\State $k$ \gets \hspace{2pt} $\sqrt{(n^2 - i^2)}$
\State $l$ \gets \hspace{2pt} $k^2$
\If{ {$l$ equals $n^2-i^2$}}
\State return 1
\EndIf
\EndFor
\State return 0
\EndProcedure
\end{algorithmic}
\end{algorithm}
\subsubsection{Time analysis}
\textbf{Best Case}: For the best case of the algorithm we need to consider the smallest possible value in the range of the number ,and the value is ${10^4}$ which itself takes ${6.10^3-2}$ iterations. Consider the time taken for the single iteration to be ${x}$ units .So the total time will be approximately t_\Omega \propto  ${(6.10^3-2)}.x$ and hence the complexity becomes overall closer to n.So the notation will be $\Omega($n$)$ .\\
\textbf{Average Case} : In this case the complexity is $\Theta($n$)$.\\
\textbf{Worst Case} : The worst case complexity can occur if the number is not indeed a Pythagorean number .So  the complexity will be O(${n}$).\\
\subsection{Better Approach}
A better method could be implementing \textbf{Euclid's} formula which is a general formula for generating Pythagorean numbers, given any arbitrary pair of integers m and n where $m>n>0$,  if $c^2=a^2+b^2$ then
\begin{center}
    \textbf{$c = m^2 + n^2 , b =2mn , a = m^2-n^2$}
\end{center}
 The triple generated by Euclid's formula is primitive if and only if m and n are coprime and not both odd. When both m and n are odd, then a, b, and c will be even, and the triple will not be primitive; however, dividing a, b, and c by 2 will yield a primitive triple when m and n are coprime and both odd.
\subsubsection{Algorithm for the better approach }
\begin{algorithm}
\begin{algorithmic}[1]
\Procedure{checkPythagorean}{$n$}
\For {$i \gets 1 \hspace{4pt} to \hspace{4pt} \sqrt{n} $}
\State $k$ \gets \hspace{2pt} $\sqrt{(n - i^2)}$
\State $l$ \gets \hspace{2pt} $k^2$
\If{ {$l$ equals $n-i^2$}}
\State return 1
\EndIf
\EndFor
\State return 0
\EndProcedure
\end{algorithmic}
\end{algorithm}
\subsubsection{Time Analysis}
\textbf{Best Case}: For the best case of the algorithm we need to consider the smallest possible value in the range of the number ,and the value is ${10^4}$ which itself takes only ${60}$ iterations. Consider the time taken for the single iteration to be ${x}$ units .So the total time will be approximately \begin{center}t_\Omega \propto\  ${60.${x}$$  $\end{center}and hence the complexity becomes overall closer to $\sqrt{n}$ which is 100.So the notation will be $\Omega(${\sqrt{n}}$)\ .\\
\textbf{Average Case} : In this case the complexity is $\Theta(\sqrt{n}$).\\\
\textbf{Worst Case} : The worst case complexity can occur if the number is not indeed a Pythagorean number .So  the complexity will be O($\sqrt{n}$).\\
\subsection{Optimal Approach}
We know that all Pythagorean number are multiples of Pythagorean prime which are prime number of the form of ${4n+1}$. Now to check whether the given number is Pythagorean number or not we can check if any of the prime factor of the given number is a Pythagorean prime or not . 
\subsubsection{Algorithm for the Optimal Approach}
\begin{algorithm}
\begin{algorithmic}[1]
\Procedure{checkPythagorean}{$n$}
\While{2 divides n}
\State $n \gets n/2$
\EndWhile
\For{$i \gets3 \hspace{2pt} to\hspace{2pt} \sqrt{n} \hspace{1pt}$ in step of 2}
\If{i divides n  and 4 divides (n-1)}
\State return 1
\EndIf
\While{i divides n}
\State $n \gets n/i$
\EndWhile
\EndFor
\If{$n>2$ and 4 divides (n-1)}
\State return 1
\EndIf
\State return 0
\EndProcedure
\end{algorithmic}
\end{algorithm}.
\subsubsection{Time Analysis}
\textbf{Best Case}:Best case holds for the least possible time to be taken from the number in the given range.In our case consider the smallest number say ${10^4}$ in this case the complexity will be given as \begin{center}
    ${t_\omega}$  $\propto\ ${log_2{10^4}+v}$  
\end{center}where ${v}$ is a lesser quantity than ${log_2{10^4}}$ .So the complexity is ${\omega(log_2{n})}$.This complexity will hold perfect for the number which are perfect powers of 2 that is of the form ${2^n}$.\\
\textbf{Worst Case} : The worst case complexity can occur if the number is not indeed a Pythagorean number. So  the loop iteration will run for ${log_2{k}}$ + ${\sqrt{z}}$ times such that ${z=(n/2^k)}$. Out of both ${\sqrt{n}}$ is greater and so the time will heavily be determined by this factor.So the complexity here will be ${o({\sqrt{n})}}$ .\\
\begin{figure}[h!]
\includegraphics[scale=0.3]{optgraph.jpg}
\caption{Comparing best case and worst case complexity in optimal approach}
\end{figure}
\subsection{Main Function}
\subsubsection{\textbf{V}ariables:}
\paragraph{1)}\textbf{a[]} is the given array of randomly generated natural numbers of size 1000 which is in the range from $10^4$ to $10^6$ .
\paragraph{2)}\textbf{min} is an integer to store the minimum Pythagorean number from the given array of numbers and is initialized to 0.
\paragraph{3)}\textbf{max} is an integer to store the maximum Pythagorean number from the given array of numbers and is initialized to 1000000.
\begin{algorithm}
\begin{algorithmic}
\Procedure{Main}{}
\For{$i \gets 0 \hspace{2pt} to \hspace{2pt}n $}
\If{$checkPythagorean(a[i]) equals \hspace{2pt} 1 $}
\If{$min > a[i]$}
\State $min \gets a[i]$
\EndIf
\If{$max < a[i]$}
\State $max \gets a[i]$
\EndIf
\EndIf
\EndFor
\EndProcedure
\end{algorithmic}
\end{algorithm}
\section{Experimental Results}
\includegraphics[scale=0.3]{expdaa1.jpg}
\caption{Figure 2:} Graph showing the comparisons of the experimental performance of the algorithm with the theoretical analysis .
\subsection{Complexity Analysis and Explanation}:
The overall complexity in given problem turns out to be ${o(k.\sqrt{n}/2 + k.\log(n))}$ where $k$ denotes the length of the array. In our experimental analysis,we considered an array of $k$ integers and then generated random numbers from $10^4$ to $10^6$ with the gap of $10^4$ to a specific value. The plot for the graph can be explained as follows : Consider a value from the randomly generated numbers if it is of the form of ${2^n.p}$ where $p$ is prime, Then the time taken in this case will be $k.({n + \sqrt{p}/2})$ which will be the worst case situation and the best case will be k.$\.log_2{n}}$ provided the loop terminates in the while block itself.Hence the overall complexity in the situation will lie in between those two values, that is
\begin{center}
    $k.\log_2{n}$ $<$ $t_{avg}$ $<$ ${k.\sqrt{n} +k.p}$
\end{center}
where k is the length of the array.

\section{Discussion and Future Work}
We have used the fact that all the Pythagorean numbers are multiples of Pythagorean prime , so we only checked if the smallest Pythagorean prime divides the given number . Other way could be , as we have been given array ranging from $10^4$ to $10^6$ , the largest prime factor of such number will be less than $10^3$ if such a number is not prime else we need to check the given number is prime and of the form 4n+1 . We can store all the Pythagorean primes less than $10^6$ in an array, and check the given number if divisible by any Pythagorean prime stored in this array . Only 80 such primes exist under $10^3$ , so a lot of computations will be minimized . Also one could use sieve method to save all the Pythagorean primes and then use the same method to check if the number is Pythagorean number or not . 
\section{Conclusion}
We analyzed various algorithms to find the largest and the smallest Pythagorean number and developed optimal approach to find the desired number from the randomly generated array of numbers . We also analyzed how Pythagorean prime factors were related to the problem . However , pre-computation of the problem to generate Pythagorean primes or with sieve method , the given problem may be more efficient if memory-time complexity trade-off is considered in favor of memory.  
\section{References}
\begin{enumerate}
\item IDAA432C (Design and Analysis of Algorithm) class lecture
\item Pythagorean triples reference -\\${https://en.wikipedia.org/wiki/Pythagorean_triple}$
\item Euclid's formula  reference\\
${http://demonstrations.wolfram.com/EuclidsFormula}$\\${AndPropertiesOfPythagoreanTriples/}$
\item Introduction to Algorithms by Cormen
\end{enumerate}
\end{document}               % End of document.


